\section{Scope and Review Method}\label{sec:methodology}\label{sec:background}
\subsection{Questions and eligibility}
We ask which hidden targets and observables appear (RQ1), how signal origin relates to receiver permission (RQ2), what establishes feasibility and transfer (RQ3), and which intervention addresses which observable at what cost (RQ4). An agent selects and executes actions using an LLM, tools, memory or orchestration; inherited serving and hardware channels need not be agent-specific.

Core eligibility requires (C1) LLM/agent execution, (C2) a hidden target, (C3) an execution-dependent auxiliary observable rather than ordinary output semantics alone, (C4) specified observation permissions, and (C5) an empirical inference evaluation. A controlled synthetic-state inference can satisfy C5; a proposed mechanism or timing gap alone cannot. Mixed papers qualify through an identifiable component. Table~\ref{tab:boundary-rules} supplies the decision rules.
\begin{table}[t]\centering\footnotesize
\caption{Observable-based eligibility. A hidden target alone does not make a semantic attack a side channel.}\label{tab:boundary-rules}
\begin{tabularx}{\columnwidth}{@{}p{.39\columnwidth}Y@{}}\toprule
Observed input & Decision and required distinction\\\midrule
Length, timing, cache response & Core when measured execution dependence supports hidden-target inference.\\
Own-site actions & Core attribution with site permission; no automatic access to other sites or private tool logs.\\
Reply states the secret & Direct disclosure; adjacent, not auxiliary-channel evidence.\\
Reply style/content identifies a user & Adjacent unless a separately evaluated auxiliary mechanism is established. WildAuth stays a mixed-view comparator.\\
Output-filter membership oracle & Can be core: the filter's execution, not ordinary semantic output inference, mediates the signal.\\
Privileged loop/fault logs & Separate diagnostic component when C1--C5 hold; does not increase confidentiality evidence.\\\bottomrule
\end{tabularx}\end{table}


The final coded collection contains \FinalN{} families available by October 8, 2026: \FinalCore{} core (\FinalA{} workflow-origin, \FinalB{} inherited-runtime/application) and \FinalAdjacent{} adjacent. Four core components concern authorized diagnostics. Other targets include content/property inference, provenance attribution and human/agent distinction; hiddenness alone establishes no confidentiality violation. Families are deduplicated studies; result records distinguish conditions within a study.

The auxiliary-only rule is applied consistently. Coding fingerprints and WildAuth remain adjacent because their evaluated views mix semantic and interaction features; feature importance is not a separately evaluated auxiliary-only attack \cite{ghaleb2026codingfingerprint,wang2026dialogue}. Metric-only fault localization qualifies as a diagnostic component despite operator authorization \cite{scheinert2026rca}. ZeroPrint's response differences, endpoint-content fingerprints and StepFinder's mixed view remain adjacent \cite{shao2026zeroprint,leshin2026endpoint,zhu2026stepfinder}.

\subsection{Purposive sampling and citation tracing}
This SoK uses purposive, mechanism-oriented sampling supplemented by fixed one-hop backward and forward citation tracing. It synthesizes decision-relevant evidence rather than estimating literature-wide prevalence. Acquisition combines native IEEE Xplore/ACM Digital Library exports, targeted discovery of preprints/proceedings with exact-title/DOI primary-source lookup, and tracing from 40 fixed seeds. Targeted searches covered arXiv and primary USENIX, NDSS and PMLR records. Appendix~\ref{app:search} gives executable native queries; Supplement S1 retains their exports and targeted query logs. All 2,237 returned backward reference occurrences and 24 distinct forward works were screened; forward lookup resolved 38 of 40 seeds. Versions are merged only with identity evidence.

We distinguish the acquired candidate frames from the coded collection. The prepared metadata frame contains 4,753 unique IDs with a recorded exclusion, referral or unresolved disposition. A bounded source reconciliation accounts for 109 overlapping decision rows in four nominated frames (57 follow-ups, six proposals, 22 backward candidates and 24 forward records). The map retains source-unavailable outcomes; it is not 109 independent full readings. At the source-search endpoint, unavailable full texts and unestablished empirical components are not admitted to the present core analysis. Unbound defense quantities are excluded from matched comparisons without removing their qualifying attack families.

The review endpoint is the named, version-bound 60-family collection. The saved native exports contain only 1/24 (IEEE) and 0/24 (ACM) fixed known papers. They are therefore inadequate as the sole acquisition route; targeted primary-source discovery and citation tracing are essential parts of this purposive design. Because database indexing and field eligibility are not established for all 24, these figures measure export presence rather than a calibrated recall rate. The unacquired broad-query frame and unresolved routes are retained in Supplement S1. Countercondition searches and source-deletion analysis test the specific synthesis, without converting it into an exhaustive systematic review.

\subsection{Coding, adjudication and reproducibility}
The extraction dictionary specifies receiver permissions, auxiliary versus semantic observations, result splitting, execution mode, utility/cost object and source/run binding. Source-scoped \NRR{}, unavailable \NRA{} and unverified \NRV{} remain distinct from replication. Code inspection can establish implementation behavior without authenticating the historical run.

One coding workflow used AI assistance in screening, extraction, source rechecks and writing. Disagreements were resolved against source text and recorded with locators. Supplement S1 retains the 23-case blinded AI exercise (17/21 determinate agreement), 332-dispute readbacks and exclusion audits. They are AI-assisted checks, not independent human coding or a human agreement estimate. The result records and conclusion dependencies expose the interpretations for direct inspection.

All \FinalCore{} core families have a scoped intervention disposition. The Supplementary Evidence supplies the 60-family index, 33 result records, 66 utility/cost field positions, source versions/locators, claim links and a fit-free verifier. Section~\ref{sec:scope} gives a complete record-to-judgment example. Supplement S6 tests source deletion and unresolved bindings; S8 supplies the pairwise and checklist audit with its coding rules and row-level reasons.
